\documentclass[../main.tex]{subfiles}
\begin{document}

\section{Đặt vấn đề}
Điểm danh là hoạt động bắt buộc và diễn ra thường xuyên tại các trường đại học nhằm quản lý chuyên cần của sinh viên. Tuy nhiên, các phương pháp điểm danh truyền thống như điểm danh bằng giấy hay gọi tên bộc lộ nhiều hạn chế: tốn thời gian của cả giảng viên và sinh viên, khó tổng hợp dữ liệu, và đặc biệt là dễ bị gian lận thông qua hình thức điểm danh hộ.

\textit{[TODO: trình bày số liệu/dẫn chứng về tình trạng điểm danh hộ, nhu cầu số hóa; nêu lý do chọn nền tảng Zalo Mini App -- tận dụng hệ sinh thái sẵn có, không cần cài app riêng.]}

\section{Mục tiêu và phạm vi đề tài}
\subsection{Mục tiêu}
Xây dựng ứng dụng điểm danh thông minh \textbf{Zimo Checkin} trên nền tảng Zalo Mini App, đáp ứng các mục tiêu:
\begin{itemize}
    \item Số hóa toàn bộ quy trình điểm danh, giảm thời gian và công sức thủ công.
    \item Chống gian lận điểm danh hộ bằng cơ chế xác thực đa lớp.
    \item Cung cấp công cụ quản lý, thống kê và phát hiện gian lận cho giảng viên và quản trị viên.
\end{itemize}

\subsection{Phạm vi}
\textit{[TODO: giới hạn phạm vi -- triển khai thử nghiệm trong môi trường đồ án; các vai trò sinh viên/giảng viên/quản trị viên; chưa tích hợp với hệ thống đào tạo của trường.]}

\section{Đối tượng sử dụng}
Hệ thống phục vụ ba nhóm đối tượng chính: \textbf{sinh viên} (thực hiện điểm danh), \textbf{giảng viên} (tạo và quản lý buổi điểm danh), và \textbf{quản trị viên} (quản lý lớp, người dùng, thống kê toàn hệ thống).

\section{Định hướng giải pháp}
\textit{[TODO: tóm tắt giải pháp -- quy trình điểm danh 4 bước, mã QR xoay HMAC, nhận diện khuôn mặt, xác minh chéo, điểm tin cậy.]}

\section{Bố cục đồ án}
Đồ án được tổ chức thành sáu chương: Chương 1 giới thiệu đề tài; Chương 2 khảo sát và phân tích yêu cầu; Chương 3 trình bày nền tảng lý thuyết và công nghệ sử dụng; Chương 4 phân tích thiết kế, triển khai và đánh giá hệ thống; Chương 5 trình bày các giải pháp và đóng góp nổi bật; Chương 6 kết luận và hướng phát triển.
\end{document}
